% 行星（综述）
% license CCBYSA3
% type Wiki

本文根据 CC-BY-SA 协议转载翻译自维基百科\href{https://en.wikipedia.org/wiki/Planet}{相关文章}。

\begin{figure}[ht]
\centering
\includegraphics[width=8cm]{./figures/c14ce525d8184ebd.png}
\caption{太阳系八大行星的尺寸比例图（自上而下、从左到右）：土星、木星、天王星、海王星（外行星），地球、金星、火星和水星（内行星）。} \label{fig_Planet_1}
\end{figure}
行星是一种巨大且呈圆形的天体，一般要求其绕恒星、恒星遗迹或褐矮星运行，并且其自身不属于这些天体。按该术语最严格的定义，太阳系共有八颗行星：类地行星（水星、金星、地球和火星）以及巨行星（木星、土星、天王星和海王星）。关于行星形成的最佳理论是星云假说，该假说认为一团星际云从星云中坍缩，形成一颗年轻的原恒星，而这颗原恒星被一个原行星盘所环绕。行星在该盘内通过由引力驱动的物质逐渐累积而成长，这一过程称为吸积。

“行星”（planet）一词源于希腊语 πλανήται（planḗtai），意为“流浪者”。在古代，这个词指的是太阳、月亮，以及在恒星背景上肉眼可见并会移动的五个光点——水星、金星、火星、木星和土星。行星在历史上具有宗教关联：多种文化将天体与神祇对应，而这种与神话和民间传说的联系，也延续到了对新发现的太阳系天体的命名方式之中。随着16至17世纪日心说取代地心说，地球本身也被认作一颗行星。

随着望远镜的发展，“行星”一词的含义被扩展，纳入了只能借助观测工具才能看见的天体：地球之外行星的卫星；冰巨行星天王星与海王星；谷神星及后来被确认属于小行星带的其他天体；以及冥王星，后来被发现是柯伊伯带中冰质天体族群中最大的成员。对柯伊伯带中其他大型天体的发现，尤其是厄里斯，引发了关于如何精确定义“行星”的争论。2006年，国际天文学联合会（IAU）通过了太阳系中“行星”的定义，将四颗类地行星与四颗巨行星归入行星范畴；而谷神星、冥王星与厄里斯则被归入“矮行星”类别。尽管如此，许多行星科学家仍持续以更宽泛的方式使用“行星”一词，包括矮行星以及如月球般呈圆形的卫星。

天文学的进一步发展使得人们发现了超过5900颗太阳系外的行星，即系外行星。这些系外行星常呈现太阳系行星所不具备的特殊性质，例如“热木星”——如飞马座51b般绕其母恒星极近距离运行的巨行星——以及极端偏心的轨道，如HD~20782~b。褐矮星与比木星更大的行星的发现也推动了对于定义边界的讨论，即如何划定行星与恒星之间的分界线。已有多颗系外行星被发现运行在其恒星的宜居带内（即行星表面可能存在液态水的区域），但地球仍是唯一已知能够承载生命的行星。宜居带的形成受多种因素影响，例如行星的大气组成（若其具有大气）。例如火星在赤道区域的白天温度足以形成液态水，只是其气压过低而无法实现。
\subsection{形成}
\begin{figure}[ht]
\centering
\includegraphics[width=8cm]{./figures/edc16278f48809b3.png}
\caption{原行星盘} \label{fig_Planet_2}
\end{figure}
\begin{figure}[ht]
\centering
\includegraphics[width=8cm]{./figures/645312c0cec1ae7e.png}
\caption{行星形成过程中原行星之间的碰撞} \label{fig_Planet_4}
\end{figure}
行星的形成机制目前尚未被完全确定。主流理论认为，行星在星云塌缩形成一层薄的气体和尘埃盘的过程中逐渐聚合而成。在中心区域形成原恒星，其周围则是一个旋转的原行星盘。通过吸积（即通过“粘性碰撞”的过程），盘中的尘埃粒子逐渐聚集质量，形成越来越大的天体。在局部区域中会出现称为微行星的质量集中体，它们通过自身的引力吸引更多物质，从而加速吸积过程。这些质量集中体变得越来越致密，最终在自引力作用下向内塌缩，形成原行星 [6]。

当一个原行星的质量增长到略高于火星的质量时，它会开始吸积一层延展的大气层 [7]，从而通过大气阻力效应 [8][9] 极大提高其对微行星的捕获效率。根据固体与气体的吸积历史不同，最终可能形成巨行星、冰巨行星或类地行星 [10][11][12]。人们认为木星、土星和天王星的规则卫星可能以类似方式形成 [13][14]；然而，海王星的卫星海卫一很可能是被捕获的 [15]。地球的月球 [16] 与冥王星的卫星卡戎 [17] 则可能是在巨型碰撞中形成。

当原恒星增大至能够点燃并成为真正的恒星时，残余的吸积盘会从内向外被光致蒸发、太阳风、Poynting–Robertson 阻力及其他效应逐渐清除 [18][19]。此后，系统中可能仍存在许多原行星，它们相互环绕或环绕恒星运行，但随着时间推移，其中许多会发生碰撞，或合并成更大的原行星，或为其他原行星提供可吸积的碎片 [20]。那些质量足够大的天体会吸尽其轨道邻域的大部分物质而成为行星。避免了碰撞的原行星可能通过引力捕获成为行星的天然卫星，或者留存在某些轨道带中，最终成为矮行星或小天体 [21][22]。
\begin{figure}[ht]
\centering
\includegraphics[width=8cm]{./figures/624858ff17708327.png}
\caption{超新星遗迹喷 ejecta 产生行星形成物质} \label{fig_Planet_5}
\end{figure}
较小的微行星的高能撞击（以及放射性衰变）会加热正在成长的行星，使其至少部分熔化。行星内部开始按照密度分异，密度更高的物质下沉至核心。\,[23] 较小的类地行星在这种聚积过程中会失去其大部分大气，但失去的气体可通过地幔逸气以及随后彗星的撞击所补充\,[24]（更小的行星会通过各种逃逸机制而失去它们获得的任何大气\,[25]）。

随着对太阳以外恒星的行星系统的发现与观测，人们得以对上述理论加以扩展、修正甚至替代。金属丰度——一个天文学术语，用来描述原子序数大于 2（即氦）的化学元素的丰度——似乎决定着恒星拥有行星的可能性。\,[26][27] 因此，与金属贫乏的第二星族（Population II）恒星相比，金属丰富的第一星族（Population I）恒星更可能拥有规模可观的行星系统。\,[28]





